% SPDX-License-Identifier: AGPL-3.0-or-later
% © Concepts 1996–2026 Miroslav Šotek. All rights reserved.
% © Code 2020–2026 Miroslav Šotek. All rights reserved.
% ORCID: 0009-0009-3560-0851 | Contact: www.anulum.li | protoscience@anulum.li
% SCPN Phase Orchestrator — arXiv manuscript of the matched-false-alarm study.
% Canonical source of the text: docs/studies/early_warning_matched_false_alarm.md
% (every number binds to the hash-sealed artefacts under examples/real_data/).

\documentclass[11pt]{article}

\usepackage[T1]{fontenc}
\usepackage[utf8]{inputenc}
\usepackage{lmodern}
\usepackage{microtype}
\usepackage{amsmath}
\usepackage{booktabs}
\usepackage[margin=1.05in]{geometry}
\usepackage[hidelinks]{hyperref}
\usepackage{url}
\urlstyle{same}
\emergencystretch=2em

\newcommand{\code}[1]{\texttt{#1}}

\title{Generic early-warning detection is at chance across five modalities;
a domain-specific detector clears the bar where the signature is
deterministic}
\author{Miroslav \v{S}otek\\
  ANULUM / Fortis Studio\\
  ORCID 0009-0009-3560-0851 \quad \href{mailto:protoscience@anulum.li}{protoscience@anulum.li}}
\date{2026}

\begin{document}
\maketitle

\begin{center}
\begin{minipage}{0.92\textwidth}
\itshape
A same-protocol replication across brain, heart, power-grid and
palaeoclimate data, with an AR(1)-Kendall-$\tau$ head-to-head against
Dakos et al.\ 2008, a dynamical-network-biomarker extension to
single-cell and bulk cancer/injury transcriptomics, and a
domain-specific head-to-head in which a power-grid modal-growth detector
beats the whole generic suite on a real corpus.
\end{minipage}
\end{center}

\begin{abstract}
Generic early-warning signals (EWS) --- the rising variance and lag-one
autocorrelation of \emph{critical slowing down}, and related
synchronisation and ordinal-entropy indicators --- are widely reported to
precede abrupt transitions in the brain, the heart, power systems and the
climate. Most of that evidence rests on a \textbf{retrospective,
per-record} test: a rising trend measured over one pre-transition window
and compared to surrogates of the \emph{same} record. We ask a different,
operational question: \textbf{at a fixed false-alarm budget, does an
early-warning detector fire on transitions more often than on
no-transition controls?} We build one domain-adaptable detector suite and
one matched-false-alarm evaluation harness, apply them unchanged (through
a per-domain adapter) to four independent labelled corpora ---
scalp-EEG seizures, cardiac atrial-fibrillation onsets, power-grid
growing oscillations, and palaeoclimate abrupt transitions --- and score
every result with a label-permutation significance test. Across all four
domains, \textbf{no detector reaches significance} (every best-member
$p \ge 0.05$): the sparse detection observed is what the matched
false-alarm rate produces by chance. Running the canonical literature
detector --- the Dakos et al.\ 2008 AR(1)-Kendall-$\tau$ trend --- through
the \textbf{same} protocol on the \textbf{same} segments confirms it: on
its own palaeoclimate records it leads 0 of 6 transitions ($p = 1.00$),
and it beats chance in none of the four domains, though on scalp EEG it
is the strongest signal anywhere (3 of 6, $p = 0.067$). We conclude that
the operational bar --- matched false alarm plus a significance test ---
is stricter than the retrospective per-record test the literature uses,
that generic EWS \emph{detection} behaves as a commodity at this bar, and
that the defensible deliverable is not a lead but the \textbf{auditable,
hash-sealed, byte-reproducible evidence} --- including the sealed
silences --- that the protocol produces. Extending the same honesty to a
\textbf{molecular fifth modality} --- dynamical-network-biomarker (DNB)
early warning of cell-fate and disease transitions, where the statistic
is cross-sample rather than sliding-window --- reaches the same
conclusion: the celebrated single-cell transition index of Mojtahedi et
al.\ 2016 clears a matched operating point on only 1 of 3 leukaemic
lineages (permutation $p = 0.27$), and the canonical GSE2565
phosgene-injury DNB benchmark does \textbf{not} beat a
\emph{selection-controlled} surrogate null (surrogate-rank $p = 0.39$)
that re-selects the biomarker module on each shuffled surrogate ---
showing the apparent DNB rise is largely a selection artefact. The same
protocol also separates commodity detection from genuine skill. Where a
domain carries a \emph{deterministic, detectable} instability signature,
a \textbf{domain-specific} detector clears the operational bar
decisively: on the PSML power grid, a modal envelope-growth detector ---
the exponential growth rate $\sigma$ of the most unstable bus's
cross-bus voltage deviation, the canonical wide-area-monitoring
instability quantity --- leads 36 of 90 growing-instability transitions
(permutation $p = 0.0001$) where every generic member is at chance (best
13 of 90, $p = 0.18$), on the identical non-circular disturbance-type
split at a matched false alarm, and a held-out half confirms it (24 of
45, $p = 0.0002$). Two pre-registered cross-dataset legs (real ISO-NE
captures and the WECC 240-bus OSL contest corpus) then show the certified
\textbf{numeric operating point does not port in either direction} ---
$\sim$32\,\% ambient false alarm on ISO-NE, zero crossings anywhere on
WECC --- so the deployable unit is the frozen shape plus per-system
matched-false-alarm calibration, and forced oscillations that switch on
effectively instantaneously offer no early-warning window at these
record lengths. A domain-specific detector does not automatically win: a
spectral detector on the murky preictal scalp EEG is itself at chance.
It is the matched-false-alarm moat that certifies which --- commodity
where the signature is absent, genuinely skilled where it is present. We
state the limitations plainly: the corpora are small (3--90 transitions;
the bulk DNB case has one exposed arm), so the tests have low power, and
``at chance'' bounds the \emph{demonstrated} skill rather than proving
early warning impossible.
\end{abstract}

\section{Background and question}\label{sec:background}

The theory of critical transitions predicts that a dynamical system
approaching a bifurcation recovers more slowly from perturbations, which
shows up as a rising variance and a rising lag-one autocorrelation of an
observable --- \emph{critical slowing down} (Scheffer et al.\ 2009). The
same idea motivates rising-synchronisation and
ordinal-transition-entropy indicators. These generic early-warning
signals have been reported before epileptic seizures, cardiac
arrhythmias, power-system instability, ecological collapse and abrupt
climate change.

Almost all of that evidence answers a \textbf{retrospective, per-record}
question: \emph{within one pre-transition window, is there a
statistically significant rising trend of the indicator, relative to
surrogate time series of that same record?} (e.g.\ the Kendall-$\tau$
trend test of Dakos et al.\ 2008). This question is valuable but it is
not the question an operator or clinician faces. The operational
question is: \emph{if I set an alarm threshold that fires on at most a
fixed fraction of no-transition situations, does the detector then fire
on genuine transitions more often than that fraction?} --- a
\textbf{matched-false-alarm} question, followed by \emph{is the
transition hit-rate above chance?}

This study builds the machinery to answer the operational question
honestly, applies it across four independent physical domains with one
shared suite, and runs the canonical literature detector through the
identical protocol as a same-segments head-to-head.

\section{Methods}\label{sec:methods}

\subsection{The detector suite}

One suite of three passive detectors reads a neutral observable bundle
and emits per-window scores:

\begin{itemize}
  \item \textbf{Critical slowing down} --- the larger of the robust
    (median/MAD) z-scores of the windowed variance and lag-one
    autocorrelation of the observable, against a leading baseline.
  \item \textbf{Rising synchronisation} --- the robust z-score of the
    Kuramoto order parameter of a population of phase oscillators.
  \item \textbf{Ordinal-transition entropy} --- the robust z-score of the
    (negated) ordinal-pattern transition entropy of the phase field.
  \item \textbf{Weighted fusion} --- the weighted mean of the three
    members.
\end{itemize}

The single-series domains (palaeoclimate) carry one scalar observable,
so only critical slowing down applies; the multi-node domains (EEG, ECG,
grid) carry a population of oscillators and use all four.

\subsection{The matched-false-alarm harness}

For each domain: a per-domain adapter turns the raw recording into the
neutral observable bundle. Each transition is scored on a \textbf{fixed
pre-onset segment ending at the annotated onset}, so every window is
pre-onset and any alarm is a genuine lead. A \textbf{no-transition null}
--- a stable, transition-free stretch --- is cut into non-overlapping
trials of the same length. Each detector's alarm threshold is set
\textbf{continuously} to the tightest value holding the trial
false-alarm rate at or below a target of 10\,\% (the quantile of the
null alarm scores, with no grid ceiling, so a variance-heavy detector is
never silently clipped). A transition is \emph{led} when the detector
alarms within its pre-onset segment. The achieved false-alarm rate is
recorded alongside every threshold.

\subsection{The permutation significance test}

A sealed detection count leaves one question open: does the detector
lead transitions \emph{more than the matched false-alarm rate explains}?
We answer it with a \textbf{label-permutation (exchangeability) test}.
Each segment --- transition or null --- alarms or not at the calibrated
threshold. Under the null hypothesis that transition segments are
indistinguishable from nulls, which of the pooled segments carry the
``transition'' label is exchangeable. We draw 10\,000 random
transition-sized subsets of the pooled alarm outcomes (fixed seed, so
the p-value is byte-reproducible), build the null distribution of the
lead count, and report the one-sided p-value with an add-one correction.
A small p-value means the detector beats the matched false-alarm rate.

\subsection{The competitor head-to-head}

The canonical literature detector is Dakos et al.\ 2008: the
\textbf{Kendall rank correlation $\tau$ of the windowed lag-one
autocorrelation against time} --- a rising-AR(1) trend. We implement it
as a per-segment score, calibrate a matched-false-alarm $\tau$ threshold
on the null segments exactly as above, and run the \textbf{same}
permutation test on the resulting alarms. Because the two detectors read
the same one-dimensional signal (the detrended proxy residual for
palaeoclimate, the cross-node Kuramoto order parameter for the
multi-node domains), the head-to-head is a same-segments, same-budget,
same-test comparison in which only the detector differs.

\subsection{Corpora (citation-only)}

\begin{table}[h!]
\centering
\small
\setlength{\tabcolsep}{4pt}
\begin{tabular}{llcl}
\toprule
Domain & Corpus & Transitions & Null \\
\midrule
Brain (scalp EEG) & CHB-MIT, chb01 (Shoeb 2009, PhysioNet) & 6 seizures & interictal records \\
Heart (ECG) & MIT-BIH AFDB (Moody \& Mark 1983) & 6 AF onsets & sinus stretches \\
Grid (PMU) & PSML 23-bus (Zheng et al.\ 2021) & 12 growing osc. & damped disturbances \\
Palaeoclimate & Dakos et al.\ 2008 records & 6 of 8 evaluated & stable pre-approach \\
\bottomrule
\end{tabular}
\caption{The four physical corpora.}
\end{table}

All raw data are public and cited; none is redistributed. Only derived,
hash-sealed, byte-reproducible evidence records are committed. Each
detector's alarm --- or silence --- is sealed into a content-addressed
(canonical-JSON SHA-256) \code{EarlyWarningEvidence} record, so the
result, including every sealed silence, is auditable and reproduces
bit-for-bit from a fresh run.

\subsection{The molecular fifth modality (DNB)}

Dynamical-network-biomarker early warning (Chen et al.\ 2012) is a
\emph{different modality}: few timepoints, many genes, and a
\textbf{cross-sample} statistic computed over the sample population at
each timepoint, so the sliding-window suite does not apply. Its index
rises to a peak at the tipping point. Two published index forms are
implemented and each is run through a modality-appropriate honest test
that keeps the same philosophy --- a controlled null and a
permutation/rank p-value:

\begin{itemize}
  \item \textbf{Single-cell (Mojtahedi et al.\ 2016).} The transition
    index is the ratio of mean signed gene--gene to mean signed
    cell--cell correlation over the curated panel (which \emph{is} the
    critical module, so no selection). We take the per-lineage index
    trajectory and its bootstrap standard error \textbf{from the
    published supplement} (not re-derived from the confounded raw qPCR);
    the pre-transition rising limb (Days 0, 1, 3) is scored by its
    least-squares slope; the matched-false-alarm null is a
    \textbf{temporally-shuffled surrogate} resampled from the published
    mean and SE; and the three lineages' alarms are tested by the same
    label-permutation core as the physical domains.
  \item \textbf{Bulk (GSE2565 phosgene lung injury).} Here the module
    must be \emph{selected}, which introduces \textbf{selection
    freedom}: the composite index
    $\mathrm{SD_{in}} \cdot \mathrm{PCC_{in}} / \mathrm{PCC_{out}}$ is
    evaluated on a module chosen to peak at the transition, so a rise is
    partly guaranteed by construction. The only honest null
    \textbf{re-selects the entire module on each surrogate}: it shuffles
    the timepoint labels across the arm's samples and reruns module
    selection from scratch, so the surrogate has the same selection
    freedom as the real analysis. The observed rising slope is ranked
    against these selection-controlled surrogates for a one-sided
    p-value, and also thresholded to a matched false alarm.
\end{itemize}

Both DNB paths seal a hash-addressed, byte-reproducible record exactly
as the physical domains do.

\subsection{Domain-specific detectors}

The generic suite reads a modality-neutral observable bundle, which is
what makes the four-domain replication possible --- but it is also why
it is a commodity. A \textbf{domain-specific} detector instead reads the
raw physical quantity whose growth \emph{is} the instability, and is
tested through the \emph{same} matched-false-alarm and permutation moat,
so its p-value is directly comparable to the generic suite on the same
segments. Two are built, one for a domain with a deterministic signature
and one for a murky domain, precisely to show the moat certifies the
difference rather than the detector always winning:

\begin{itemize}
  \item \textbf{Power-grid modal envelope-growth.} When a disturbance
    leaves an electromechanical mode under-damped, the amplitude of the
    cross-bus voltage deviation grows exponentially --- the real part
    $\sigma$ of the mode's eigenvalue is positive, the canonical
    wide-area-monitoring early-warning quantity (Kundur 1994). The
    detector estimates $\sigma$ directly as the slope of the log
    deviation envelope. Two orthogonal, physically-motivated refinements
    are added and selected on a development half only: a \textbf{focal}
    aggregation takes the growth rate of the \emph{most unstable bus}
    (instability is localised to a mode/bus cluster, so the per-bus
    maximum, un-diluted, is the signal; the matched-false-alarm
    threshold is calibrated on that same per-bus maximum over the nulls,
    absorbing the multiplicity), and a \textbf{recency weighting}
    up-weights the later samples (a real instability accelerates toward
    the disturbance). Every variant tried is disclosed in
    \S\ref{sec:modal}.
  \item \textbf{Scalp-EEG spectral rise.} The seizure-prediction
    literature reports a rising high-frequency to low-frequency
    band-power ratio before onset (e.g.\ Mormann et al.\ 2007). The
    detector takes the beta(13--30\,Hz)/delta(0.5--4\,Hz) band-power
    ratio per channel, its Kendall-$\tau$ rising trend over the
    pre-onset segment, with the same whole-head and focal aggregations.
    Its bands are fixed a-priori from the literature, not tuned.
\end{itemize}

The grid labels are \textbf{non-circular}: a transition is any
generator-trip scenario and a null any damped bus-fault or branch-trip
scenario --- the label is the disturbance \emph{type}, a physical
annotation independent of the growth statistic the detector measures, so
a growing-oscillation detector is never scored against a growth-defined
label. Three PSML scenarios carry a non-monotonic time column (a
non-physical, negative inferred sampling rate) and are dropped, the
count recorded in the sealed payload. The operating point (focal
aggregation, recency weighting) is chosen on the even-index development
half and validated on the odd-index held-out half before the
full-corpus comparison, so no operating point is tuned to the reported
answer.

\section{Results}\label{sec:results}

\subsection{SCPN suite across four domains}\label{sec:fourdomain}

At a matched false-alarm rate (target 10\,\%), detection is sparse
everywhere, and \textbf{no detector reaches significance in any domain}:

\begin{table}[h!]
\centering
\begin{tabular}{llcc}
\toprule
Domain & Best member & Led / N & Permutation $p$ \\
\midrule
Brain (EEG) & critical slowing down & 2 / 6 & 0.215 \\
Heart (ECG) & synchronisation & 2 / 6 & 0.219 \\
Grid (PMU) & critical slowing down & 3 / 12 & 0.195 \\
Palaeoclimate & critical slowing down & 1 / 6 & 0.514 \\
\bottomrule
\end{tabular}
\caption{Best generic suite member per domain at a matched 10\,\% false alarm.}
\end{table}

The fusion never leads more transitions than its best single member, so
there is no robust fusion advantage. The observed lead counts are within
what the matched false-alarm rate produces by chance (expected counts
0.6--1.7 across the domains).

\subsection{Head-to-head: SCPN vs Dakos AR(1)-Kendall-$\tau$}\label{sec:headtohead}

Running the canonical Dakos detector through the identical protocol on
the identical segments:

\begin{table}[h!]
\centering
\begin{tabular}{lcc}
\toprule
Domain & SCPN best member (led/N, $p$) & Dakos AR(1)-Kendall-$\tau$ (led/N, $p$) \\
\midrule
Palaeoclimate & CSD 1/6, $p = 0.514$ & 0/6, $\mathbf{p = 1.000}$ \\
Grid & CSD 3/12, $p = 0.195$ & 1/12, $\mathbf{p = 0.713}$ \\
Heart & sync 2/6, $p = 0.219$ & 0/6, $\mathbf{p = 1.000}$ \\
Brain (EEG) & CSD 2/6, $p = 0.215$ & \textbf{3/6, $p = 0.067$} \\
\bottomrule
\end{tabular}
\caption{Same-segments, same-budget, same-test head-to-head.}
\end{table}

On its own palaeoclimate records --- the data Dakos et al.\ analysed ---
the AR(1)-Kendall-$\tau$ detector, at a matched operating point, leads
\textbf{zero} of six transitions ($p = 1.00$): it does not beat chance,
and it does no better than the SCPN suite. The same holds for grid and
heart. The one exception is scalp EEG, where the rising-AR(1) trend is
the \textbf{strongest signal anywhere in the study} (3 of 6,
$p = 0.067$) --- marginally short of significance, and better than the
SCPN detectors on that domain.

\subsection{The honest reading}

The result is not ``nothing works''. It is: \textbf{at a matched
operating point and a permutation significance test, nothing reaches
significance at these corpus sizes} --- neither the SCPN suite nor the
canonical literature detector, in any of the four domains. The
literature's positive EWS results come from a \emph{retrospective
per-record} significance test that this operational protocol is stricter
than. The one place a detector comes close --- the AR(1) trend on scalp
EEG --- is consistent with the seizure-prediction literature's use of
autocorrelation features and warrants a larger, higher-power EEG corpus.

\subsection{The molecular fifth domain (DNB)}\label{sec:dnb}

The same conclusion holds when the honesty is carried to a molecular
modality:

\begin{table}[h!]
\centering
\footnotesize
\setlength{\tabcolsep}{4pt}
\begin{tabular}{llll}
\toprule
Case & Test & Result & $p$ \\
\midrule
Single-cell (Mojtahedi 2016) & matched-FA + permutation, 3 lineages & 1 / 3 lineages clear 10\,\% & 0.266 \\
Bulk (GSE2565 phosgene) & selection-controlled surrogate rank & rise does not beat surrogates & 0.385 \\
\bottomrule
\end{tabular}
\caption{The molecular modality under modality-appropriate honest nulls.}
\end{table}

The single-cell transition index rises 2.8-fold toward the erythroid
bifurcation --- an unmistakable effect size --- yet at its published
four-timepoint resolution only the strongest of three lineages clears a
matched false-alarm operating point, and the corpus-level permutation
test is not significant ($p = 0.27$); the binding constraint is the
temporal resolution, not the effect size. The bulk GSE2565 case is
sharper: the phosgene-exposed arm's apparent DNB rise (slope 1.15) sits
\emph{below} the 90th percentile of surrogates that were allowed to
\textbf{re-select the biomarker module on shuffled time} (surrogate-rank
$p = 0.385$), and it barely exceeds the air control's own
selected-module rise ($p = 0.414$). Once the null is granted the same
selection freedom as the analysis, the celebrated bulk DNB rise is
largely a \textbf{selection artefact}. Both molecular cases are
hash-sealed and byte-reproducible.

\subsection{Domain-specific detectors: a real signature beats the bar}\label{sec:modal}

The generic result is that \emph{modality-neutral} detection is a
commodity. The complementary result is that a \emph{domain-specific}
detector, tested through the identical moat, beats it decisively --- but
only where the domain carries a deterministic signature.

On the PSML power grid, the non-circular disturbance-type split gives
\textbf{90 growing-instability (generator-trip) transitions and 88
damped (bus-fault / branch-trip) nulls} --- a larger and
label-independent corpus than the growth-selected 12 of
\S\ref{sec:fourdomain}. Both the modal detector and the generic suite
read the \emph{same} two-second pre-onset segments of the \emph{same}
scenarios, calibrated to the \emph{same} 10\,\% matched false alarm on
the \emph{same} nulls (achieved 9.1\,\%); only the detector differs. At
its validated operating point (focal aggregation, recency weighting):

\begin{table}[h!]
\centering
\begin{tabular}{lcc}
\toprule
Detector & Transitions led / 90 & Permutation $p$ \\
\midrule
generic critical slowing down & 13 / 90 & 0.185 \\
generic weighted ensemble & 11 / 90 & 0.324 \\
generic ordinal-transition entropy & 8 / 90 & 0.629 \\
generic synchronisation & 3 / 90 & 0.976 \\
\textbf{modal envelope-growth (focal, recency)} & \textbf{36 / 90} & \textbf{0.0001} \\
\bottomrule
\end{tabular}
\caption{Grid modal head-to-head at the identical matched operating point.}
\end{table}

The domain-specific detector leads \textbf{36 of 90} transitions
(40\,\%) at $p = 0.0001$, beating every generic member --- the best of
which (critical slowing down, 13/90, $p = 0.19$) is at chance. The
\textbf{held-out half} confirms it: 24 of 45 transitions (53\,\%) at
$p = 0.0002$, the unbiased estimate, sealed in the artefact.

The winning operating point was reached by an explicit, disclosed
variant search on the development half (choice made on development,
reported on held-out --- no head-to-head data touched during selection):

\begin{table}[h!]
\centering
\small
\begin{tabular}{lcc}
\toprule
Variant & Development led / 45, $p$ & Held-out led / 45, $p$ \\
\midrule
whole-network log-slope $\sigma$ & 11, 0.051 & 14, 0.009 \\
per-bus maximum $\sigma$ (focal) & 14, 0.008 & 22, 0.0001 \\
matrix-pencil, PCA-dominant mode & 10, 0.072 & 3, 0.80 \\
matrix-pencil, per-bus maximum & 8, 0.189 & 6, 0.39 \\
Hilbert-envelope $\sigma$, per-bus & 2, 0.90 & 3, 0.80 \\
\textbf{recency-weighted focal $\sigma$} & \textbf{16, 0.002} & \textbf{24, 0.0002} \\
\bottomrule
\end{tabular}
\caption{The disclosed variant search behind the winning operating point.}
\end{table}

The gold-standard wide-area-monitoring estimator --- \textbf{matrix-pencil
modal damping} --- recovers a planted $\sigma$ exactly on synthetic data,
yet on the short two-second, 238\,Hz windows it is at chance (held-out
$p = 0.39$--$0.80$); the robust direct log-envelope slope wins on real
noisy segments. A recency-ratio sensitivity sweep (2--8) beats the
unweighted per-bus detector across the whole range, so the choice is
robust, not knife-edge.

The counterpoint confirms this is not detector-worship. The
\textbf{scalp-EEG spectral} detector --- a beta/delta band-power rise,
a-priori bands, tested through the same moat on the chb01 seizures ---
is itself at chance (1 of 6 led, $p \approx 0.56$, in both whole-head
and focal aggregations): the preictal state is murky and the corpus
small, so a domain-specific detector on a domain without a clean
signature fares no better than the generic suite. \textbf{The moat is
what certifies the difference} between a genuine, physically-grounded
early warning (grid modal growth) and a plausible but empty one
(preictal spectral rise) --- the same honest test that flags the
commodity detectors as at chance flags the domain-specific one as
genuinely skilled.

\subsection{The streaming operating point: offline certification is
necessary, not sufficient}\label{sec:stream}

The grid head-to-head of \S\ref{sec:modal} is a \emph{per-window} result
--- it scores the pre-onset window. Run the same certified detector as a
\textbf{live stream} (scoring every sliding window causally, as a PMU
feed would deliver it) and the operating point is markedly stricter,
because a damped bus-fault or branch-trip produces a \emph{short
transient growth} window the continuous monitor alarms on. At the
certified per-window threshold the stream leads 82 of 90 transitions but
false-alarms on \textbf{73\,\%} of damped scenarios --- operationally
useless. Recalibrating for the stream at a matched stream-level false
alarm, with a persistence debounce and an exponential-fit-quality gate
that rejects a fault's step-like transient (a fault fits an exponential
poorly, a genuine instability well), the fit-quality-gated detector
leads \textbf{11 of 45 held-out transitions (24\,\%)} at a matched
10\,\% stream false alarm --- well below the offline per-window 40\,\%.
The gap is physical: distinguishing a sustained instability from a
damped fault \emph{online} requires observing the damping sign over
several cycles, a lead-time / discrimination limit the pre-onset-window
benchmark does not face. This sharpens the moat's lesson once more: even
a genuinely skilled per-window detector must be re-certified for its
deployment modality, and the honest streaming operating point ---
hash-sealed in
\path{examples/real_data/psml_modal_growth/grid_modal_stream_operating_point.json}
--- is the auditable deliverable, not the flattering per-window figure.
The final step from a live alarm to a \emph{decision} is taken honestly
too: a passive, review-only advisory
(\path{assurance.grid_early_warning_advisory}) turns each stream
alarm into a claim-bounded operator record that surfaces $\sigma$, the
most-unstable bus, and the certified operating point, and seals the
$\sim$24\,\% recall as a first-class field --- so the advisory never
actuates, is a reason to inspect rather than a guarantee, and makes
explicit that the \emph{absence} of an advisory is not evidence of
stability.

\subsection{The grid moat's winning form does not transfer to EEG: a
model-mismatch boundary}\label{sec:eegtransfer}

A natural next question is whether the grid detector's \emph{winning
form} --- the exponential envelope-growth rate, focal aggregation, and
the fit-quality gate that lifts its streaming operating point
(\S\ref{sec:stream}) --- carries to another domain. We test it on the
murky scalp EEG: the same a-priori seizure feature as \S\ref{sec:modal}
(the per-channel beta/delta band-power ratio trajectory) is scored not
by a Kendall-$\tau$ rank trend but by the \textbf{exponential growth
rate of its envelope}, with the identical focal aggregation, recency
weighting and $R^2$-gate the grid uses, through the same
matched-false-alarm + permutation moat on the same six chb01 preictal
segments. It does \textbf{not} transfer: the exponential-growth score
leads \textbf{0 of 6} preictal segments ($p = 1.00$) in both
aggregations, and the fit-quality gate --- built to reject a fault's
\emph{non-exponential} step-like transient --- rejects the preictal
\emph{signal} too (every gated score is clamped, the achieved false
alarm falls to zero, 0 of 6 led), below even the a-priori rank trend's
1 of 6. A robustness grid over the analysis window (512/1024/2048
samples), recency weighting (1.0/3.0) and gate (off/0.5) leaves the
modal focal score at 0 of 6 at \textbf{every} setting, so this is a
model mismatch, not a parameter artefact. The reason is physical: the
grid statistic presupposes a \emph{linear instability's eigenvalue} ---
the deviation grows as a true exponential, which the fit and its quality
gate are built to measure --- whereas the preictal band-power rise, if
present, is a slow, noisy, non-exponential drift the rank trend reads
better than an exponential fit. The exponential-growth moat is therefore
\textbf{grid-specific}, and this boundary is itself hash-sealed in
\path{examples/real_data/chb01_seizures/seizure_modal_transfer.json}.
(A shorter 600\,s horizon does surface a spectral rank-trend lead of 3
of 6, $p = 0.04$; with one subject, $n = 6$, and a single horizon
forking path that does not survive multiplicity, it is disclosed in the
artefact as exploratory and explicitly \textbf{not} claimed --- the same
discipline that keeps the programme's positive results honest applies to
its near-misses.)

\subsection{The moat cannot even be posed on DNB: a resolution-limit
boundary}\label{sec:dnbtransfer}

The second candidate for the grid moat's transfer is the molecular DNB
modality of \S\ref{sec:dnb}, and it fails for a different, structural
reason. The DNB early-warning signal is a genuine
critical-slowing-down \emph{rise} of the transition index toward the
tipping point --- but it is measured over only a handful of
pre-transition timepoints: the single-cell Mojtahedi rising limb is
\textbf{three} points (Days 0, 1, 3, from the published Table~S2), and
the bulk GSE2565 exposed-arm limb is \textbf{four} (0.5, 1, 4, 8\,h). On
three-to-four points the grid statistic's discriminating machinery has
no purchase: a fit-quality gate needs enough of a trajectory to tell an
exponential from any other monotone rise, and on so few points
\emph{every} rise fits an exponential well enough ($R^2 = 0.67$--$0.86$
across all four trajectories) to pass a 0.5 gate --- so the gate keeps
every one and rejects none, exactly the discrimination it is meant to
provide, gone. The exponential growth rate, moreover, simply re-orders
the three single-cell lineages in the same order as the existing linear
slope ($0.61 > 0.48 > 0.22$ vs $0.128 > 0.092 > 0.028$), adding no
separating information. So the modal-growth transfer collapses to the
linear slope the existing DNB detector already uses, and the moat's
winning form cannot even be \emph{posed}. This boundary is hash-sealed
in \path{examples/real_data/mojtahedi_fate/dnb_modal_transfer.json}.

Together, \S\ref{sec:eegtransfer} and \S\ref{sec:dnbtransfer} bound the
moat from both sides. The grid detector transfers only to a domain that
carries \textbf{both} a genuine exponential instability (a linear mode's
eigenvalue, which fits an exponential and whose fault-vs-instability
contrast a fit-quality gate can read) \textbf{and} a trajectory long
enough to fit and gate that form. The power grid has both; scalp EEG has
a resolved trajectory but a non-exponential preictal rise the gate
rejects (a model mismatch, \S\ref{sec:eegtransfer}); DNB has a genuine
early-warning rise but only a three-to-four-point limb the gate cannot
resolve (a resolution limit, \S\ref{sec:dnbtransfer}). The
exponential-growth moat is therefore not a general early-warning method
but a \textbf{domain-specific} one whose applicability is itself an
auditable, sealed result --- the same honest boundary-drawing that
separates the commodity generic detectors from the genuinely skilled
grid detector.

\subsection{Cross-system external validation against eigenvalue ground
truth}\label{sec:andes}

The PSML result is on one 23-bus corpus, and its held-out ceiling
($\approx 50\,\%$) proved robust across method classes (matrix-pencil
modal identification) and aggregations (per-bus focal, dominant spatial
mode) --- none beat the ceiling at a matched false alarm, and their
union of led transitions is unharvestable there because combining raises
the null floor. To test the detector on \emph{independent} systems, we
run a stronger, non-circular check: on the IEEE 39-bus New England
system and the Kundur two-area system (both in the open-source ANDES
simulator), we sweep the operating point, take the true dominant
electromechanical mode's growth rate $\sigma$ from ANDES
\textbf{small-signal eigenvalue analysis} --- a completely different
method from the detector's time-domain envelope slope --- and score a
ringdown from a small disturbance at each point. The detector's
\textbf{coherent} aggregation (the cross-bus mean, or the dominant
spatial mode) recovers the true $\sigma$ trend on both systems (Spearman
$\rho = 0.87$ on IEEE-39, 0.60--0.66 on Kundur), so the quantity the
detector estimates --- the dominant mode's damping --- generalises
across systems and simulators. The \textbf{focal} aggregation, the PSML
winner, does \textbf{not} transfer ($\rho = -0.59$ and $-0.27$): on a
slow coherent inter-area mode the per-bus maximum locks onto spurious
local excursions rather than the network mode. The growth-rate quantity
is therefore universal, but the best \textbf{aggregation is
regime-dependent} --- focal for PSML's fast, localised oscillations,
coherent for the slow inter-area modes here --- a new, actionable
finding hash-sealed in
\path{examples/real_data/grid_external_validation/grid_eigenvalue_external_validation.json}.
The test is honest about its limits: a simulated (not field-PMU) ground
truth, stable operating points only (the well-damped benchmarks do not
enter a clean growing-oscillation regime, so this is a damping-ranking
test, not a stable-versus-unstable classification), and a fourteen-point
sweep per system.

\subsection{Critical-slowing-down external validation against the
analytic eigenvalue}\label{sec:csd}

The grid section validates the \emph{oscillatory} detector against a
simulated eigenvalue. The single-series \textbf{critical-slowing-down}
detector --- the one the palaeoclimate capstone runs on the Dakos
records --- admits an even harder, non-circular test, because its ground
truth can be written in closed form rather than simulated. We sweep a
control parameter quasi-statically toward a codimension-one bifurcation,
integrate the stochastic normal form at each operating point, and read
the shipped detector's two indicators --- the lag-one autocorrelation
and the variance. The true recovery rate $\lambda$ is the normal-form
\textbf{Jacobian eigenvalue} analytically ($-2\sqrt{\mu}$ for the fold /
saddle-node, $\mu$ for the supercritical pitchfork). On both independent
classes --- one with a quadratic nonlinearity, one with a globally
confining cubic one --- the detector's autocorrelation channel recovers
the true $\lambda$ (Spearman $\rho = 0.98$ fold, 0.96 pitchfork), and
because the lag-one autocorrelation of a linear-response process is
$\exp(\lambda \, \Delta t)$, the implied rate $\ln(\mathrm{AR1})/\Delta t$
tracks $\lambda$ \textbf{in magnitude}, not merely in rank (Pearson
$\rho = 0.98$, 0.96) --- a stronger statement than the rank-only grid
test can make. The variance channel rises in step ($\rho = 0.98$ on
both) as the stationary variance $\sigma^2/2|\lambda|$ diverges. So the
quantity the detector estimates on the Dakos records --- a falling
recovery rate --- is exactly what a bifurcation presents, confirmed
against a first-principles reference and hash-sealed in
\path{examples/real_data/csd_bifurcation/csd_bifurcation_external_validation.json}.
The limits are stated plainly: a quasi-static per-point sweep (an
independent stationary run at each control value, not one
non-stationary approach), additive noise only, and scalar reduced normal
forms rather than the full high-dimensional systems.

\subsection{The Hopf bridge: which family recovers the eigenvalue on an
oscillatory mode}\label{sec:hopf}

The grid section (\S\ref{sec:andes}) validates the \emph{oscillatory}
envelope-growth family on real systems; \S\ref{sec:csd} validates the
\emph{single-series} autocorrelation family on \emph{non-oscillatory}
bifurcations. A \textbf{Hopf} bifurcation is where the two meet: it
presents an \emph{oscillatory} critical slowing down (eigenvalue
$\alpha \pm i\omega$ with $\alpha \to 0^{-}$), so both families point at
the same mode against one analytic ground truth. We sweep the Hopf
parameter $\alpha$, integrate the stochastic normal form
($dr = (\alpha r - r^3)\,dt$ on the amplitude, $d\theta = \omega\,dt$ on
the phase), and read each family with its shipped detector. The
\textbf{envelope-growth family recovers $\alpha$ in rank \emph{and}
magnitude} (Spearman $\rho = 0.97$, mean $|\sigma - \alpha| = 0.04$ at
an adequate ringdown SNR), and the recovery is
\textbf{frequency-invariant} (identical across 0.2--1.2\,Hz, because it
fits the envelope, not the oscillation). The \textbf{autocorrelation
family tracks $\alpha$ but not its magnitude} (mean
$|\mathrm{rate} - \alpha| = 0.69$): its lag-one autocorrelation is
confounded by the oscillation, pinned near $\cos(\omega \, \Delta t)$,
so $\ln(\mathrm{AR1})/\Delta t$ does not estimate $\alpha$. This
completes a \textbf{regime map}: the eigenvalue's real part is the
universal quantity, but the magnitude-correct estimator is
regime-dependent --- envelope-growth for an oscillatory mode (here, and
the grid), the autocorrelation for a non-oscillatory one
(\S\ref{sec:csd}). Hash-sealed in
\path{examples/real_data/hopf_bridge/hopf_bridge_external_validation.json}.
The one honest sensitivity is stated as a curve, not hidden: the
envelope magnitude recovery \textbf{degrades with ringdown SNR} ($\rho$
falls from 0.97 to 0.62 as the ringdown noise rises), a physical floor
--- a decay cannot be read below the noise it sinks into.

\subsection{The eigenvalue regime map, consolidated}\label{sec:regimemap}

The three external validations above
(\S\ref{sec:andes}--\S\ref{sec:hopf}) are one result seen from three
angles. Each pins a shipped detector against an \emph{independent}
ground truth for the mode's eigenvalue --- one simulated (ANDES
small-signal analysis), two analytic (normal-form Jacobian) --- and
together they map which estimator recovers that eigenvalue in which
regime:

\begin{table}[h!]
\centering
\footnotesize
\setlength{\tabcolsep}{3.5pt}
\begin{tabular}{lllll}
\toprule
Mode & Systems & Eigenvalue truth & Estimator & Recovery \\
\midrule
Oscillatory, real (\S\ref{sec:andes}) & IEEE-39, Kundur & ANDES $\sigma$ & envelope, coherent & rank ($\rho$ 0.87 / 0.60--0.66) \\
Non-oscillatory (\S\ref{sec:csd}) & fold, pitchfork & normal-form $\lambda$ & autocorrelation & \textbf{magnitude} ($\rho$ 0.98 / 0.96) \\
Oscillatory (\S\ref{sec:hopf}) & Hopf normal form & normal-form $\alpha$ & envelope $\sigma$ & \textbf{magnitude} ($\rho$ 0.97) \\
\bottomrule
\end{tabular}
\caption{The eigenvalue regime map across three independent ground truths.}
\end{table}

The \textbf{eigenvalue's real part --- the growth rate, or its
sign-flipped damping --- is the universal quantity}, validated across
three independent ground-truth methods. The \textbf{magnitude-correct
estimator is regime-dependent}, and the split is not arbitrary but
mechanical: on an \emph{oscillatory} mode the lag-one autocorrelation is
confounded by the oscillation (pinned near $\cos(\omega \, \Delta t)$,
so $\ln(\mathrm{AR1})/\Delta t \ne \alpha$), while the envelope carries
the decay cleanly --- so envelope-growth wins on the grid and the Hopf
mode. On a \emph{non-oscillatory} mode there is no oscillation for the
envelope to fit, and the linear-response autocorrelation is exactly
$\exp(\lambda \, \Delta t)$ --- so the autocorrelation channel wins on
the fold and the pitchfork. Pick the estimator by whether the mode
oscillates; the quantity you are estimating is the same eigenvalue in
every case.

\subsection{The regime map extends to the collective coordinate}\label{sec:kuramoto}

\S\ref{sec:andes}--\S\ref{sec:regimemap} validate the eigenvalue on
\emph{scalar} systems (normal forms) and on a per-bus grid mode. The
last test carries it to the model SCPN is built around, one dimension
harder: the eigenvalue now lives on the \textbf{emergent collective
coordinate} of a high-dimensional system, not a one-line normal form.
For the noisy Kuramoto model of $N = 512$ phase oscillators (Lorentzian
frequencies, half-width $\gamma$, phase diffusion $D$), the incoherent
state loses stability at the mean-field critical coupling
$K_c = 2(\gamma + D)$ --- the $K_c = 2\gamma$ of the Ott--Antonsen
reduction, generalised for noise --- and just below it the fundamental
mean-field mode is real and non-oscillatory with the analytic eigenvalue
$\lambda(K) = (K - K_c)/2$ (Sakaguchi 1988). The regime map's
non-oscillatory prescription therefore predicts the
\emph{autocorrelation} family should recover $\lambda$; sweeping $K$
below $K_c$ and reading the shipped detector on the mean field
$Z = (1/N) \sum \exp(i\theta)$ confirms it --- \textbf{but only through
the right observable}. Both the signed real part $\mathrm{Re}(Z)$ and
the order-parameter amplitude $|Z|$ track $\lambda$ in \textbf{rank}
(Spearman $\rho = 0.97$ and 0.98: the collective coordinate slows down
as $K \to K_c$), yet \textbf{only the signed $\mathrm{Re}(Z)$ recovers
$\lambda$ in magnitude}, fitting it with slope 1.19 ($\approx 1$)
because its lag-one autocorrelation is $\exp(\lambda \, \Delta t)$; the
folded amplitude $|Z|$ a practitioner usually watches fits with slope
2.66 --- it ranks the eigenvalue but cannot size it. So the map's
magnitude tier survives the jump from a scalar normal form to a
512-oscillator order parameter, with one added, actionable rule: to read
the distance to the synchronisation threshold, monitor the signed
mean-field component, not the order-parameter magnitude, and use a long
enough autocorrelation window that the finite-window bias does not
steepen the estimate. Hash-sealed in
\path{examples/real_data/kuramoto_synchronization/kuramoto_synchronization_external_validation.json}.
Honest limits: the eigenvalue is the $N \to \infty$ mean-field result
while the runs are finite-$N$ (an $O(1/\sqrt{N})$ correction, which is
why the signed slope is near, not exactly, one), a quasi-static
per-coupling sweep, a Lorentzian law with additive noise, and coupling
below $K_c$ only.

\subsection{The collective map has two regimes, mirroring the whole
study}\label{sec:bimodal}

\S\ref{sec:kuramoto} recovered the eigenvalue on a \emph{non-oscillatory}
collective transition. A \textbf{bimodal} frequency law is the
oscillatory counterpart, and completes a two-regime map inside the
Kuramoto model --- the collective analogue of the Hopf bridge
(\S\ref{sec:hopf}). For a symmetric bimodal Lorentzian (peaks at
$\pm\omega_0$, half-width $\Delta$), the incoherent state loses
stability at $K_c = 4\Delta$, and for $\omega_0 > K/4$ the fundamental
eigenvalue is \textbf{complex} (Martens et al.\ 2009, via the
Ott--Antonsen reduction):
$\lambda_{\pm}(K) = K/4 - \Delta \pm \sqrt{(K/4)^2 - \omega_0^2}$, so
$\mathrm{Re}(\lambda) = K/4 - \Delta$ and
$\Omega = \sqrt{\omega_0^2 - (K/4)^2}$. Below $K_c$ the order parameter
is a damped oscillation --- exactly the regime where the Hopf bridge
found the \emph{envelope} family, not the autocorrelation, sizes the
eigenvalue. Ringing down from a partially-coherent start ($N = 4000$,
$\omega_0 = 1.5$, $\Delta = 0.5$) and reading two families confirms it:
the \textbf{envelope-growth family on the sub-population order parameter
$|Z_+|$ recovers $\mathrm{Re}(\lambda)$ in magnitude} (Spearman
$\rho = 0.99$, fitted slope 0.93 $\approx$ 1), while the
\textbf{autocorrelation family is confounded by the oscillation}
$\Omega$ (fitted slope 15.8 --- it ranks the eigenvalue but cannot size
it). The observable is the subtlety: the \emph{global} mean field is a
standing wave of two counter-rotating modes, so $|Z|$ oscillates to zero
and cannot be fit, whereas the sub-population $Z_+$ is a single complex
mode whose modulus decays smoothly as $\exp(\mathrm{Re}(\lambda)\,t)$.
The measured oscillation frequency matches the analytic $\Omega$ to a
mean absolute error of 0.11 ($\sim$7\,\%) --- the analytic $\Omega$ is
nearly constant across the sweep, so this is a value check, not a rank
one --- confirming the mode oscillates at the predicted frequency, i.e.\
the eigenvalue is genuinely complex. Hash-sealed in
\path{examples/real_data/kuramoto_bimodal/kuramoto_bimodal_hopf_external_validation.json}.
So the collective coordinate carries \textbf{both} regimes of the
study's map: non-oscillatory $\to$ autocorrelation
(\S\ref{sec:kuramoto}), oscillatory $\to$ envelope-growth (here),
exactly as the scalar normal forms (\S\ref{sec:csd}, \S\ref{sec:hopf})
and the grid (\S\ref{sec:andes}) do. Honest limits: a noiseless
finite-$N$ ringdown (the $O(1/\sqrt{N})$ floor bounds the envelope),
the oscillatory regime $\omega_0 > \Delta$ only, and coupling below
$K_c$.

\subsection{Cross-dataset portability: the operating point does not
travel, in either direction}\label{sec:portability}

The strong grid result (\S\ref{sec:modal}--\S\ref{sec:stream}) is
certified on one corpus. The External Validation Program's cross-dataset
leg (E2.G) asks the deployment question directly: does the certified
detector --- its \textbf{shape frozen}, no variant search permitted ---
survive a dataset change? Two independent corpora answer it, one real
and one synthetic, under protocols pre-registered before the first
detector run and sealed with every amendment disclosed.

\textbf{The frozen numeric operating point does not port --- measured in
both directions.} On three real ISO-NE PMU captures of documented
sustained oscillations (UTK test-cases library), the PSML per-window
threshold ($1.3203\,\sigma/\mathrm{s}$, two-second windows) fires on
$\sim$32\,\% of pre-onset ambient frequency windows --- against the
certified 9.09\,\% --- so the certified constant is far too \emph{hot}
there. On the 13 forced-oscillation cases of the 2021 IEEE-NASPI OSL
contest (WECC 240-bus synthetic PMU, bus-voltage observable --- the
certified observable family), the same frozen threshold crosses
\textbf{nothing}: 0 of 611 pooled ambient windows and 0 transition
windows --- far too \emph{conservative}. A threshold certified on
238\,Hz fault-transient voltage envelopes is a corpus-specific quantity,
which is why the product's deployment step calibrates per system on the
target system's own ambient data and re-seals.

\textbf{The frozen shape with local matched-false-alarm calibration is
honest about what remains.} On ISO-NE (window = five cycles of the
documented mode, threshold at a matched 10\,\% false alarm on pooled
pre-onset ambient windows only), the detector catches 1 of 3 events ---
the 1.13\,Hz regional mode, 57.1\,s before the estimated onset --- with
permutation $p = 0.332$ at $n = 3$: a case study, no significance claim.
On WECC, the corpus structure itself becomes the finding: with the onset
search pinned at the exact forcing start ($t = 30$\,s, known by
simulation design), 9 of 13 cases are already past three times the
ambient in-band envelope at the forcing start --- an effectively
\textbf{instantaneous onset with no early-warning window at all} --- and
of the 4 cases with a resolvable window the calibrated shape leads 2
($+5.4$\,s, $+4.1$\,s; $p = 0.41$, not significant). A secondary,
descriptive branch shows the same operating point \emph{detects} the
oscillation after onset in 12 of 13 cases with $\sim$1\,s median latency
--- but its count is chance-compatible on 60\,s regions at a 10\,\%
window false alarm (the per-case chance bound is sealed alongside), so
it is reported without any significance claim.

The two legs close the loop the moat opened: the certified quantity (the
growth rate $\sigma$) and the certified \emph{procedure}
(matched-false-alarm calibration plus sealed evidence) travel across
datasets; the certified \emph{number} does not, and forced oscillations
that switch on instantaneously offer no early-warning window for any
detector at these record lengths. Both records are hash-sealed and
byte-reproducible ---
\path{examples/real_data/iso_ne_forced_oscillation/} and
\path{examples/real_data/wecc_240_osl/} --- with pre-registration
appendices, disclosed protocol amendments, and integrity tests that pin
the honest headlines, including the negatives.

\section{Discussion}\label{sec:discussion}

\textbf{Detection is a commodity; the moat is the evidence.} Across four
independent physical domains --- and a fifth, molecular one --- generic
early-warning \emph{detection} at an honest operating point is sparse
and, by a permutation or selection-controlled test, at chance. This is
not a defect of one suite: the canonical Dakos detector fares no better
on its own data, and the celebrated single-cell and bulk DNB benchmarks
do not clear a modality-appropriate honest null either. What is
\emph{not} a commodity is the auditable, reproducible, claim-bounded
envelope the protocol produces --- a matched-false-alarm operating
point, a permutation p-value, and a hash-sealed
\code{EarlyWarningEvidence} record for every transition,
\textbf{including the sealed silences}. A positive early-warning claim
should be required to clear this operational bar; most published EWS
results have only cleared the retrospective per-record one.

\textbf{A domain-specific detector wins where the signature is
deterministic --- and the moat certifies which.} The commodity result is
about \emph{modality-neutral} detectors. When the domain carries a
physically deterministic, directly-measurable instability --- a growing
electromechanical mode, whose exponential envelope \emph{is} the
eigenvalue crossing into instability --- a detector that reads that
quantity beats the whole generic suite decisively (grid modal growth,
36/90, $p = 0.0001$, versus every generic member at chance). But the
same detector \emph{idea} on a murky domain (the preictal spectral rise)
is itself at chance. So the deliverable is not ``build domain-specific
detectors'' as a slogan --- a domain-specific detector does not
automatically win --- but the \textbf{matched-false-alarm moat that
certifies the difference}: the identical honest test that flags
commodity detection as at chance also licenses a genuinely-skilled
detector where a real signature exists, and refuses to license a
plausible-but-empty one where it does not. Value is unlocked per-domain,
and only when the physics carries a detectable signal that clears the
operational bar.

\textbf{The winning detector estimates an eigenvalue, and that is
externally checkable.} The domain-specific grid result is not a black
box that happens to correlate with trouble: the quantity it reads is the
dominant mode's growth rate --- the real part of a physical eigenvalue
crossing toward instability --- and
\S\ref{sec:andes}--\S\ref{sec:regimemap} confirm it against three
\emph{independent} ground truths (a power-system simulator's
small-signal eigenvalues, and the closed-form Jacobian eigenvalues of
two bifurcation classes). This turns the moat from a
matched-false-alarm \emph{ranking} into a claim about a named,
first-principles quantity, and it yields a small operational rule
(\S\ref{sec:regimemap}): the eigenvalue's real part is universal, but
read it with the envelope-growth family on an oscillatory mode and the
autocorrelation family on a non-oscillatory one --- because on an
oscillatory mode the lag-one autocorrelation is confounded by the
oscillation, and on a non-oscillatory one there is no envelope to fit. A
detector that clears the operational bar \emph{and} estimates a quantity
you can check against an eigenvalue is a stronger deliverable than
either property alone.

\textbf{Selection freedom is a hidden operating cost.} The bulk-DNB case
adds a specific lesson: when the detector \emph{chooses} its own module
(or feature set) to peak at the transition, an honest null must be
granted the same choice. A surrogate that re-selects from scratch on
shuffled time absorbs most of the apparent signal --- so a fair null for
a self-selecting detector is not an afterthought but the whole test.

\textbf{The evaluation gap is the finding.} The difference between
``there is a detectable rising trend in this record's approach''
(retrospective, per-record, vs surrogates) and ``the detector fires on
transitions more than on controls at a fixed false-alarm budget''
(operational, matched, significance-tested) is exactly the gap between
the literature's positive results and this study's null. That gap ---
not any single detector --- is what the four-domain replication and the
Dakos head-to-head expose.

\section{Limitations}\label{sec:limitations}

Stated plainly, because they bound the claim:

\begin{itemize}
  \item \textbf{Small corpora, low power.} With 6--12 transitions per
    physical domain --- and only 3 lineages (single-cell) or 1 exposed
    arm (bulk) in the molecular one --- the significance tests have
    limited power; ``at chance'' bounds the \emph{demonstrated} skill
    and does \textbf{not} prove early warning is impossible. A modestly
    skilled detector could remain non-significant at this $n$.
  \item \textbf{One competitor.} Only the Dakos AR(1)-Kendall-$\tau$
    detector was run head-to-head on the physical domains; modern
    approaches (e.g.\ deep-learning tipping-point predictors) were not
    tested and could behave differently.
  \item \textbf{Single-subject EEG; within-record climate null.} The EEG
    corpus is one subject (chb01); the palaeoclimate null is the stable
    pre-approach interval of each record, which carries that record's
    own variability and so is conservative.
  \item \textbf{Parameterisation.} One reasonable choice of window,
    step, baseline and target false alarm was used per domain; a sweep
    was not performed.
  \item \textbf{DNB caveats.} The single-cell index is taken from the
    published supplement (its bootstrap SE drives the surrogate), not
    re-derived from the confounded raw qPCR; the bulk case has one
    exposed arm, so it is a single-transition surrogate test, not a
    corpus, and its module search is a greedy reimplementation of the
    DNB selection rather than the authors' exact procedure. These bound
    the molecular result to ``no significance at this resolution / under
    a selection-controlled null'', not a refutation of the DNB method.
  \item \textbf{Domain-specific operating point.} The grid modal
    detector's operating point (focal aggregation, recency weighting)
    was chosen by a variant search on a development half and reported on
    a held-out half, so the held-out lead count (24/45, $p = 0.0002$) is
    the unbiased estimate and the full-corpus count (36/90) is the
    fitted-model deployment figure; both are sealed. The corpus is one
    power-system dataset at one sampling rate with two-second windows,
    so the strong grid result bounds \emph{demonstrated} skill on PSML
    growing-oscillation transitions, not a universal grid claim. The
    scalp-EEG spectral counterpoint is one subject (chb01, six
    seizures), so its at-chance result is low-power like the generic EEG
    result.
  \item \textbf{Cross-dataset legs (\S\ref{sec:portability}).} The
    ISO-NE corpus is three usable events with a thin pooled-null base
    (six windows), so its 1-of-3 lead is a case study, not a powered
    test. The WECC corpus is synthetic (TSAT time-domain simulation, not
    field PMU), its records are 90\,s, and its forced oscillations
    mostly switch on effectively instantaneously --- so its
    early-warning null result partly reflects the corpus structure (no
    pre-onset growth to detect) rather than detector skill alone; the
    detection branch is descriptive because a 10\,\% window false alarm
    on 60\,s regions alarms by chance with high probability, a units
    mismatch the sealed record discloses explicitly.
\end{itemize}

\section{Reproduction}\label{sec:reproduction}

Every result regenerates deterministically (no randomness in the
pipelines; the permutation test uses a fixed seed). With each raw corpus
in a directory:

\begin{verbatim}
python bench/early_warning_leadtime_eeg.py     DATA OUT   # brain
python bench/early_warning_leadtime_cardiac.py DATA OUT   # heart
python bench/early_warning_leadtime_grid.py    DATA OUT   # grid
python bench/early_warning_leadtime_climate.py DATA OUT   # palaeoclimate
python bench/head_to_head_ar1_kendall.py OUT \
    --climate-dir C --grid-dir G --cardiac-dir H --eeg-dir E
python -m bench.early_warning_dnb          OUT   # single-cell DNB
python -m bench.early_warning_dnb_bulk     DATA OUT   # bulk GSE2565 DNB
python -m bench.grid_modal_head_to_head    PSML OUT   # grid modal vs generic
python -m bench.early_warning_leadtime_isone --data-dir ISONE --output OUT
python -m bench.early_warning_leadtime_wecc  --data-dir WECC  --output OUT
\end{verbatim}

The single-cell DNB capstone reads only the embedded published summary,
so it needs no external data; the bulk capstone reads the citation-only
GSE2565 files; the grid modal head-to-head reads the citation-only PSML
scenarios. The sealed evidence, aggregate results (with the permutation
block), the head-to-head comparisons, and the sealed grid
modal-vs-generic result
(\path{examples/real_data/psml_modal_growth/}) are committed under
\path{examples/real_data/} in the source repository
(\url{https://github.com/anulum/scpn-phase-orchestrator}). A fresh run
reproduces every \code{content\_hash} bit-for-bit; the integrity tests
recompute each sealed hash from the committed payload alone.

\section{Data availability}\label{sec:data}

Raw corpora are public and cited, not redistributed:

\begin{itemize}
  \item CHB-MIT Scalp EEG Database ---
    \url{https://physionet.org/content/chbmit/}
  \item MIT-BIH Atrial Fibrillation Database ---
    \url{https://physionet.org/content/afdb/}
  \item PSML power-system dataset (Zheng et al.\ 2021) --- the 23-bus
    millisecond-level PMU measurements.
  \item Dakos et al.\ 2008 palaeoclimate records ---
    \url{https://github.com/earlywarningtoolbox/datasets}
  \item Mojtahedi et al.\ 2016 single-cell transition index ---
    Table~S2 of the paper (index + bootstrap SE per lineage per day).
  \item GSE2565 phosgene lung-injury expression ---
    \url{https://www.ncbi.nlm.nih.gov/geo/query/acc.cgi?acc=GSE2565}
  \item ISO-NE sustained-oscillation PMU captures --- UTK oscillation
    test-cases library,
    \url{https://web.eecs.utk.edu/~kaisun/Oscillation/actualcases.html}
    (Maslennikov et al.\ 2016).
  \item WECC 240-bus forced-oscillation cases --- 2021 IEEE-NASPI OSL
    contest via the same library,
    \url{https://web.eecs.utk.edu/~kaisun/Oscillation/contestcases.html}
    (NREL reduced 240-bus model; credit DOE/NREL/Alliance).
\end{itemize}

\begin{thebibliography}{99}

\bibitem{scheffer2009} Scheffer M, Bascompte J, Brock WA, et al.
Early-warning signals for critical transitions. \emph{Nature} 461:53
(2009).

\bibitem{dakos2008} Dakos V, Scheffer M, van Nes EH, Brovkin V,
Petoukhov V, Held H. Slowing down as an early warning signal for abrupt
climate change. \emph{PNAS} 105(38):14308 (2008).

\bibitem{dakos2012} Dakos V, Carpenter SR, Brock WA, et al. Methods for
detecting early warning signals of critical transitions in time series.
\emph{PLoS ONE} 7(7):e41010 (2012).

\bibitem{shoeb2009} Shoeb A. Application of machine learning to
epileptic seizure onset detection and treatment. PhD thesis, MIT (2009);
CHB-MIT via PhysioNet.

\bibitem{moody1983} Moody GB, Mark RG. A new method for detecting atrial
fibrillation using R-R intervals. \emph{Computers in Cardiology} 10:227
(1983).

\bibitem{goldberger2000} Goldberger AL, Amaral LAN, Glass L, et al.
PhysioBank, PhysioToolkit, and PhysioNet. \emph{Circulation}
101(23):e215 (2000).

\bibitem{zheng2021} Zheng X, et al. A multi-scale time-series dataset
with benchmark for machine learning in electricity (PSML) (2021).

\bibitem{kundur1994} Kundur P. \emph{Power System Stability and
Control.} McGraw-Hill (1994) --- small-signal (modal) stability; a
mode's eigenvalue real part is its growth rate.

\bibitem{mormann2007} Mormann F, Andrzejak RG, Elger CE, Lehnertz K.
Seizure prediction: the long and winding road. \emph{Brain} 130(2):314
(2007).

\bibitem{chen2012} Chen L, Liu R, Liu Z-P, Li M, Aihara K. Detecting
early-warning signals for sudden deterioration of complex diseases by
dynamical network biomarkers. \emph{Scientific Reports} 2:342 (2012).

\bibitem{mojtahedi2016} Mojtahedi M, Skupin A, Zhou J, et al. Cell fate
decision as high-dimensional critical state transition. \emph{PLoS
Biology} 14(12):e2000640 (2016).

\bibitem{sciuto2005} Sciuto AM, Phillips CS, Orzolek LD, Hensley JL,
Chang L-Y, Nadadur SS. Genomic analysis of murine pulmonary tissue
following carbonyl chloride inhalation. \emph{Chemical Research in
Toxicology} 18(11):1654 (2005); GSE2565.

\bibitem{maslennikov2016} Maslennikov S, Wang B, Zhang Q, Ma F, Luo X,
Sun K, Litvinov E. A test cases library for methods locating the sources
of sustained oscillations. IEEE PES General Meeting (2016).

\bibitem{yuan2020} Yuan H, Biswas RS, Tan J, Zhang Y. Developing a
reduced 240-bus WECC dynamic model for frequency response study of high
renewable integration. 2020 IEEE/PES T\&D (2020).

\bibitem{sakaguchi1988} Sakaguchi H. Cooperative phenomena in coupled
oscillator systems under external fields. \emph{Progress of Theoretical
Physics} 79(1):39 (1988).

\bibitem{martens2009} Martens EA, Barreto E, Strogatz SH, Ott E, So P,
Antonsen TM. Exact results for the Kuramoto model with a bimodal
frequency distribution. \emph{Physical Review E} 79:026204 (2009).

\end{thebibliography}

\end{document}
